\chapter{DarkVec: NLP per scoprire gli scanner di Internet}
\label{cap:darkvec}
\fonti{\rif{2026-11-23}{23/11}}{8-Darkvec, pp.\ 1--90}

DarkVec è il caso di studio con cui il corso applica Word2Vec (cap.~\ref{cap:word2vec}) a un problema di cybersecurity: capire
\textbf{chi scandisce Internet} a partire dal traffico che arriva a una darknet. Gli IP sorgente si trattano come parole, il traffico
come testo, e lo skip-gram impara un embedding per ogni IP senza alcuna label. Il riferimento principale è Gioacchini et al.,
\emph{DarkVec: automatic analysis of darknet traffic with word embeddings} (CoNEXT 2021); la slide iniziale elenca anche i lavori
successivi dello stesso gruppo (i-DarkVec, embedding con GNN temporali, embedding federati, analisi dinamica dei cluster). È il gruppo
di ricerca del docente, che però non è tra gli autori.

\begin{nota}[la domanda di ricerca (\rif{2026-11-23}{23/11})]
Nasce attorno al 2020, dopo la botnet \textbf{Mirai}, il primo attacco riuscito che ha usato dispositivi IoT vulnerabili per campagne di
denial of service su Internet. Gli scanner non hanno un nome, solo un IP, e l'IP può cambiare: bloccarlo dopo averlo visto attaccare
molte volte è già tardi. Se so che un IP è di Mirai e le sue ``compagnie'' sono sempre le stesse, posso generalizzare agli altri IP e
accorgermi prima dell'attacco?
\end{nota}

% =====================================================================
\section{Internet: uno spazio enorme e sparso}
Lo spazio IPv4 conta $2^{32} \approx 4{,}2$ miliardi di indirizzi, ciascuno con $2^{16}$ porte TCP e UDP (slide 8, pp.~6--11). In
questo spazio vivono servizi e dispositivi, utenti e attaccanti. Per attaccare Internet, o per proteggerlo, serve
\textbf{visibilità}: sapere cosa c'è e chi lo sta cercando.

\subsection{Gli scanner}
Gli \textbf{scanner} sondano Internet per scoprire quali dispositivi e servizi sono attivi, quali vulnerabilità hanno, quali sono i
percorsi di routing (slide 8, pp.~15--19). Possono essere su tutto Internet o \textbf{mirati} (per esempio su un servizio che si vuole
attaccare), \textbf{benigni} (aziende come Google o Censys che mappano Internet) o \textbf{malevoli}: ogni dispositivo vulnerabile
scoperto allarga la superficie d'attacco.
\begin{itemize}
  \item Strumenti: Nmap (1997), Masscan e ZMap (2013).
  \item Il traffico di scansione è aumentato di circa 30 volte dal 2015 al 2024.
  \item Episodi noti: la botnet \textbf{Mirai} (settembre 2016) e l'incidente \textbf{Log4j/Log4Shell} (dicembre 2021).
\end{itemize}
Per osservarli servono \textbf{punti di osservazione} (\en{vantage points}), attivi o passivi: più sono numerosi e distribuiti nel mondo,
più si vede. Il Politecnico ha darknet in sede, in altre università italiane, in Brasile, negli Stati Uniti e nel cloud (telescopi
molto più grandi esistono negli Stati Uniti). La verità di base resta quasi sempre ignota: che un IP fosse davvero di una botnet si può
solo supporre.

% =====================================================================
\section{Network telescope (darknet)}
\begin{definizione}[darknet / network telescope --- slide 8, p.~30]
Insieme di indirizzi IP che \textbf{non ospitano alcun servizio}: non inviano mai richieste e non rispondono mai. Tutto il traffico che
ricevono è quindi \textbf{non richiesto}: scansioni (malevole e benigne), \en{backscattering} delle vittime, configurazioni errate.
\end{definizione}
È una misura \textbf{passiva}: la darknet occupa IP e porte e ascolta, ma non risponde e non partecipa alla scansione. Non offrendo
servizi né pubblicizzandosi, per costruzione chiunque la contatti ha uno scopo diverso dall'usare un servizio: c'è sempre qualcosa da
imparare, se si riesce a elaborare i dati su larga scala.
\begin{integrazione}[backscattering]
Quando un attaccante usa IP sorgente falsificati (per esempio in un DoS), la vittima risponde a quegli indirizzi: se cadono in una
darknet, il telescopio vede le risposte e quindi, indirettamente, l'attacco.
\end{integrazione}
I dati raccolti da un telescopio (slide 8, pp.~35--36) hanno quattro caratteristiche che rendono impraticabile l'analisi manuale:
\begin{center}
\small
\begin{tabularx}{\textwidth}{>{\raggedright\arraybackslash}p{3.3cm} X >{\raggedright\arraybackslash}p{3.6cm}}
\toprule
\textbf{Caratteristica} & \textbf{In pratica} & \textbf{Strumento del corso}\\
\midrule
Volume enorme & un telescopio /24 (256 indirizzi) riceve più di un milione di pacchetti al giorno & ---\\
Dati categorici & il traffico è fatto soprattutto di attributi categorici (IP, porte, protocolli) & embedding\\
Niente label & non c'è una corrispondenza tra IP e utenti reali, tanto meno per gli attaccanti & self-supervised learning\\
Alta dinamicità & le attività nascono e cambiano di continuo & ---\\
\bottomrule
\end{tabularx}
\end{center}
Il docente la chiama la parte ``a pensiero veloce'' dell'AI: dati massicci e mutevoli che vanno trattati subito, dove un transformer o un
modello linguistico sarebbe troppo lento. Poiché la darknet non risponde, dopo la scansione non arriva alcuna richiesta: si vedono solo IP
sorgente e porta cercata, cioè dati \textbf{categorici}. Le label mancano, salvo i pochi IP già segnalati. Il monitoraggio manuale si
riduce a filtri sul volume (fermarsi quando il traffico cresce troppo), che intervengono quando ormai è tardi.

L'obiettivo è passare dal traffico grezzo, categorico e senza label a una \textbf{rappresentazione numerica ricca}: un vettore per ogni
IP sorgente, prodotto da un embedder self-supervised (slide 8, pp.~38--45). Sugli embedding si costruiscono poi i \textbf{downstream
task}: trovare attività coordinate, anomaly detection, cercare gli IP più vicini (se la rete li ha messi insieme, hanno qualcosa in
comune), classificare se ci sono abbastanza dati. La parte difficile è il \emph{come}: trasformare solo IP e porte in qualcosa di molto
più ricco.

% =====================================================================
\section{L'idea di DarkVec}
\begin{definizione}[DarkVec --- slide 8, pp.~48--49]
Word2Vec impara una rappresentazione del \textbf{significato di una parola} dal contesto in cui compare. DarkVec impara una
rappresentazione del \textbf{comportamento di uno scanner} (un IP sorgente) dalle sue ``compagnie''.
\textbf{Ipotesi}: gli scanner che compaiono insieme nel tempo (\textbf{co-occorrenza temporale}) si comportano in modo simile nello
stesso periodo.
\end{definizione}
L'intuizione viene dalla natura di una botnet: tanti host che arrivano \textbf{insieme} a fare cose simili, sullo stesso servizio
oppure su porte o vantage point diversi ma nello stesso momento. Riprendendo Firth: si capisce cosa vuole un IP dalle compagnie che
frequenta (\rif{2026-11-23}{23/11}). In una darknet:
\begin{itemize}
  \item chi arriva \textbf{per errore} non compare con costanza insieme ad altri;
  \item gli scanner \textbf{benigni} di solito sono noti (le grandi aziende pubblicano i propri IP);
  \item chi compare con costanza insieme ad altri e non è in elenco restringe il cerchio: perché è lì?
\end{itemize}

\subsection{Il traffico come testo}
\dacompletare{Questa sezione e le seguenti seguono le slide (pp.~50--90); il docente le spiega nella lezione del 24/11,
da integrare quando sarà disponibile la registrazione.}

Il traffico raccolto in un mese è una tabella con timestamp, IP sorgente, porta di destinazione e protocollo (slide 8, pp.~50--62).
DarkVec lo modella come un insieme di documenti:
\begin{itemize}
  \item gli \textbf{IP sorgente} sono le \textbf{parole};
  \item la \textbf{sequenza degli IP} in ordine di arrivo è una \textbf{frase}.
\end{itemize}
Una sola frase con tutto il traffico però mette vicini IP che non hanno niente in comune: un IP che sonda la porta 22 (SSH) può
arrivare subito dopo uno che sonda la 445 (SMB). Le \textbf{porte contano}: per questo si costruisce \textbf{una frase per servizio}
(porta/protocollo). Così IP1 e IP4, che sondano entrambi la 22/TCP, finiscono nello stesso contesto; IP2 e IP3, sulla 445/TCP, in un
altro. Si scartano gli IP con meno di 10 pacchetti al giorno.

\subsection{La pipeline}
Raccolta dei dati $\rightarrow$ preprocessing (frasi di IP per servizio) $\rightarrow$ \textbf{DarkVec} (skip-gram, self-supervised)
$\rightarrow$ \textbf{task finale} non supervisionato (slide 8, pp.~64--69). Il docente fa notare che è la stessa struttura delle
pipeline dei lab. Nell'articolo ogni IP diventa un punto in uno spazio a $V$ dimensioni: le coppie (input, predizione) dello
skip-gram si costruiscono dentro una finestra $[-c, +c]$ attorno a ogni IP, separatamente per ogni servizio.

% =====================================================================
\section{Il task finale: clustering degli scanner}
Nell'articolo (slide 8, pp.~72--73): DarkVec addestrato su un mese di traffico (circa 100\,000 IP, 486 milioni di skip-gram), poi
\textbf{clustering} degli embedding dell'ultimo giorno con l'algoritmo di \textbf{Louvain} (rilevamento di comunità su grafi). Due
scenari:
\begin{itemize}
  \item \textbf{A}: gli IP degli scanner sono noti (es.\ quelli di un progetto di scansione come Censys): i cluster mostrano gruppi di
        IP con attività regolari e coordinate nel tempo;
  \item \textbf{B}: gli IP non sono noti: compaiono di continuo IP nuovi.
\end{itemize}

% =====================================================================
\section{Il limite di Word2Vec e i-DarkVec}
Il problema delle parole nuove (cap.~\ref{cap:word2vec}) qui è centrale (slide 8, p.~76):
\begin{itemize}
  \item Word2Vec impara solo le parole viste in training: non sa rappresentare parole (IP) nuove, e di scanner nuovi se ne vedono
        ogni giorno;
  \item nel linguaggio il significato di una parola cambia raramente, mentre uno scanner può cambiare comportamento in fretta;
  \item non si può riaddestrare da zero ogni volta che compare qualcosa di nuovo.
\end{itemize}
\textbf{i-DarkVec} passa da un addestramento unico a uno \textbf{incrementale} (slide 8, pp.~77--83):
\begin{center}
\small
\begin{tabularx}{\textwidth}{X X}
\toprule
\textbf{DarkVec} & \textbf{i-DarkVec}\\
\midrule
documenti con tutto il traffico osservato & traffico diviso in \textbf{batch giornalieri}, documenti per ogni giorno\\
embedding imparati in un solo passaggio (un mese di dati) & embedding aggiornati giorno per giorno, usando quelli già noti come
\textbf{ancore}\\
--- & si continua ad aggiornare quando arriva traffico nuovo, sfruttando quanto già imparato, senza ripartire da zero\\
\bottomrule
\end{tabularx}
\end{center}
Secondo la slide i risultati sono di fatto gli stessi che si otterrebbero con DarkVec.

% =====================================================================
\section{Conclusioni del docente}
DarkVec è un esempio di ciò di cui tratta il corso (slide 8, p.~89): conoscere il problema reale (esperti di cybersecurity), imparare
tecniche e buone pratiche di AI, applicarle al problema. Domande aperte:
\begin{itemize}
  \item cosa succede con gli IP nuovi (per esempio quelli del mese dopo)? $\rightarrow$ i-DarkVec;
  \item si può trasferire la conoscenza da una rete all'altra (embedding imparati sulla darknet del Politecnico usati su una darknet
        brasiliana)? $\rightarrow$ embedding federati e incrementali.
\end{itemize}
\begin{esame}[possibile domanda d'esame (slide 8, p.~89)]
DarkVec impara dalla \textbf{co-occorrenza temporale} degli scanner. Funzionerebbe sul traffico di tutti i giorni?
\end{esame}
\begin{integrazione}[uno spunto di risposta]
In una darknet ogni pacchetto è non richiesto e gli scanner coordinati arrivano insieme sugli stessi servizi, quindi stare vicini nel
tempo dice qualcosa sul comportamento. Nel traffico ordinario la vicinanza temporale è dominata da utenti e servizi legittimi che non
hanno nulla in comune: l'ipotesi va verificata, e probabilmente servono contesti diversi (per esempio sessioni dello stesso host).
\end{integrazione}
Una demo è disponibile sul GitHub del docente (\code{MatteoBoffa/DarkVec\_Demo}).

\begin{sintesi}
\begin{itemize}
  \item Darknet = IP senza servizi: ricevono solo traffico non richiesto. Dati enormi, categorici, senza label, molto dinamici.
  \item DarkVec = Word2Vec sul traffico: IP = parola, sequenza di IP per servizio = frase; ipotesi della co-occorrenza temporale.
  \item Skip-gram $\rightarrow$ un embedding per IP $\rightarrow$ clustering (Louvain) per trovare gruppi di scanner coordinati.
  \item Limite: IP nuovi fuori vocabolario e comportamenti che cambiano $\rightarrow$ i-DarkVec, embedding aggiornati ogni giorno.
\end{itemize}
\end{sintesi}
